% -*- TeX:RNW:UK -*-
\documentclass[10pt]{beamer}
\usetheme{metropolis}
%\useinnertheme{rectangles}
\setbeamercovered{%
still covered={\opaqueness<1->{15}},
again covered={\opaqueness<1->{40}}}

\hypersetup{colorlinks,linkcolor=black,urlcolor=brown,citecolor=brown}

\usepackage{amsmath,amssymb,amsthm}
\usepackage{unicode-math}

% We set the Lucida OTF fonts as default
\usepackage{fontspec}
\setmainfont{Lucida Bright OT}
\setsansfont{Lucida Sans OT}
\setmonofont{Lucida Console DK}[Scale=MatchLowercase]

\newfontfamily\typiconsfont{Typicons}[Scale=0.9]
\newcommand\typicons[1]{{\typiconsfont\symbol{#1}}}
\newcommand\Attention{\colorbox{yellow}{\textcolor{black}{\typicons{"E136}}}\xspace}
\newcommand\Advanced{\colorbox{black}{\textcolor{white}{\typicons{"E04E}}}\xspace}

\newfontfamily\webglyphsfont{WebHostingHub-Glyphs}[Scale=0.7]
\newcommand\webglyphs[1]{{\webglyphsfont\symbol{#1}}}
\newcommand\Discussion{\colorbox{white}{\textcolor{black}{\webglyphs{"F134}}}\xspace}
\newcommand\DExamples{\colorbox{black}{\textcolor{white}{\webglyphs{"F134} examples?}}}
\newcommand\Reading{\colorbox{black}{\textcolor{white}{\webglyphs{"F0C1}}}\xspace}

\newfontfamily\meteoconsfont{Meteocons}
\newcommand\meteocons[1]{{\meteoconsfont\symbol{#1}}}
\newcommand\meteosun{\meteocons{"0042}}
\newcommand\meteosuncloud{\meteocons{"0048}}
\newcommand\meteorain{\meteocons{"0052}}
\newcommand\meteosolidsun{\meteocons{"0031}}

\newfontfamily\KRfarmfont{KR Back On The Farm}
\newcommand\KRfarm[1]{{\KRfarmfont\symbol{#1}}}
\newcommand\farmplant{\KRfarm{"0049}}

\usepackage{polyglossia}
\setdefaultlanguage[variant = british, ordinalmonthday = false]{english}

\usepackage[style=authoryear-comp,firstinits,sortcites,maxcitenames=2,%
    mincitenames=1,maxbibnames=10,minbibnames=10,uniquename=mininit,%
    uniquelist=minyear,sortfirstinits=true]{biblatex}
\addbibresource{../references/ecophys.bib}
%\renewcommand{\bibfont}{\small}

\usepackage{abbrev}
\DeclareGraphicsExtensions{.png,.jpg,.jpeg,.pdf}

\begin{document}

\title{PBIO-141\\Sensory and Physiological Ecology of  Plants}
\subtitle{11: Carbon Assimilation (B)}
\author{Pedro J. Aphalo}
\date{January-February 2018}
\institute[Univ.\ of Helsinki]{M.Sc.\ in Plant Biology, University of Helsinki\\[2ex] \url{http://blogs.helsinki.fi/aphalo/}}


  \begin{frame}
    \maketitle
  \end{frame}

  \begin{frame}[c]
    \begin{center}
      \begin{small}
        \copyright 2006-2018 by Pedro J. Aphalo\\
       University of Helsinki, Finland.\\
        \textcolor{blue}{\url{http://blogs.helsinki.fi/aphalo/}}\\[2ex]
      \end{small}

      \begin{footnotesize}
        Sensory and Physiological Ecology of  Plants: `11: Carbon Assimilation (B)' by Pedro J. Aphalo is licensed under a Creative Commons Attribution-ShareAlike 4.0 International License.

        Illustrations and text quoted from copyrighted sources is excluded from this license and their use should respect the original licenses.\\[2ex]
      \end{footnotesize}

      \includegraphics[width=6em]{../figures/copyright/by-sa}
    \end{center}
  \end{frame}

  \begin{frame}
    \frametitle{Outline}
    \tableofcontents
  \end{frame}

\section{Biochemistry: C4}

\begin{frame}{C4 plants}
    \centering
    \includegraphics[width=0.85\textwidth]{figures/CornField.eps}\\
    {\small Maize is a C4 grass. Maize field in Liechtenstein (from \url{http://en.wikipedia.org/wiki/Maize}).}
\end{frame}

\begin{frame}{Discovery}
    The C4 pathway was discovered by M. D. Hatch and C. R. Slack, two
    Australian researchers, in 1966, so it is sometimes called the
    Hatch-Slack pathway.
\end{frame}

%\begin{frame}{C3 and C4 leaves}
%    \includegraphics[width=\textwidth]{figures/Gurevitch2.10CDC3C4leaves.eps}\\
%    {\small The anatomy of typical C4 leaves is different to that of C3 leaves
%    \autocite[from][]{GurevitchEtAl2006}. However C4 metabolism is also possible within
%    individual cells \autocite[for details see][latest edition only]{TaiZei2006}.}
%\end{frame}
%
\begin{frame}{C4 anatomy}
    \centering
    \includegraphics[height=0.65\textheight]{figures/Heldt8.8C4anatomy.eps}\\
    {\small C4 or ``kranz'' anatomy \autocite[from][]{Heldt1997}.}
\end{frame}

\begin{frame}{C3 vs.\ C4}
    \centering
    \includegraphics[height=0.75\textheight]{figures/Heldt8.6C3C4_CO2.eps}\\
    {\small Overall differences \autocite[from][]{Heldt1997}.}
\end{frame}

\begin{frame}{C4 mechanism}
   \begin{itemize}
        \item The ``principle'' of C4 metabolism is that the
        concentration of \COtwo\ at the place where Rubisco is
        located is increased by ``pumping'' it either to special
        cells or exceptionally to a particular region of a single
        cell.
        \item The ``pumping'' is a metabolic pathway that includes
        carboxylation, transport and decarboxylation.
        \item Assimilation takes place in the same way as in C3
        plants: in the Calvin cycle.
        \item However there is a previous carboxylation step, that
        yields a molecule with four carbon atoms as product (that is
        the origin of the name C4).
    \end{itemize}
\end{frame}

\begin{frame}{C4 mechanism}
    \centering
    \includegraphics[width=\textwidth]{figures/Heldt8.9C4principle.eps}\\
    {\small C4 ``principle'' \autocite[from][]{Heldt1997}.}
\end{frame}

\begin{frame}{Types of C4 metabolism}
    The different types differ in the cell compartment where
    decarboxylation takes place, and the enzyme involved.
    (Carboxylation is catalyzed by PEP-carboxylase in all types.)
    \begin{itemize}
        \item NADP-malic enzyme type (maize, sugarcane).
        \item NAD-malic enzyme type.
        \item PEP-carboxykinase enzyme type.
    \end{itemize}
\end{frame}

%\begin{frame}[toc=,bm=]{C4 mechanism}
%%    \vspace{-0.04\textheight}
%    \twocolumn[lcolwidth=0.57\linewidth,rcolwidth=0.37\linewidth]{
%    \includegraphics[height=\textheight]{figures/Heldt8.11C4mechanismMalic.eps}}{%
%    \vspace{0.8\textheight}
%    {\small NADP-malic enzyme type\\
%    \autocite[from][]{Heldt1997}.}}
%\end{frame}
%
%\begin{frame}[toc=,bm=]{C4 mechanism}
%%    \vspace{-0.04\textheight}
%    \twocolumn[lcolwidth=0.67\linewidth,rcolwidth=0.27\linewidth]{
%    \includegraphics[height=\textheight]{figures/Heldt8.14C4mechanismNADMalic.eps}}{%
%    \raggedright \vspace{0.3\textheight}
%    {\small NAD-malic enzyme type \autocite[from][]{Heldt1997}.}}
%\end{frame}
%
%\begin{frame}[toc=,bm=]{C4 mechanism}
%%    \vspace{-0.04\textheight}
%    \twocolumn[lcolwidth=0.67\linewidth,rcolwidth=0.27\linewidth]{
%    \includegraphics[height=\textheight]{figures/Heldt8.15C4mechanismPEPcarboxykinase.eps}}{%
%    \raggedright \vspace{0.3\textheight}
%    {\small PEP-carboxykinase enzyme type \autocite[from][]{Heldt1997}.}}
%\end{frame}
%
%\begin{frame}[toc=,bm=]{C4 mechanism}
%%    \vspace{-0.04\textheight}
%    \includegraphics[width=\textwidth]{figures/Heldt8.15C4mechanismPEPcarboxykinaseLegend.eps}\\
%    {\small PEP-carboxykinase enzyme type \autocite[from][]{Heldt1997}.}
%\end{frame}

\begin{frame}{Calvin cycle}
    The reactions of the Calvin cycle occur some in mesophyll
    chloroplasts and some in the bundle sheath chloroplasts, with
    metabolites shuttling between these two types of cells.
\end{frame}

\begin{frame}{C4 in one cell}
    Recently is has been discovered that C4 metabolism is also
    possible within a single cell. The two ends of an elongated cell fulfill
    the roles of ``mesophyll'' and ``bundle sheath'' respectively in \species{Borszczowia
    aralocaspica} (Chenopodiaceae). In another species without elongated cells, there
    are concentric regions.
\end{frame}

\section{Biochemistry: CAM}

\begin{frame}{CAM plants}
    \centering
    \includegraphics[height=0.75\textheight]{figures/CAMphoto-1.eps}
    {\small Some arid habitat plants, typically with CAM.}
\end{frame}

\begin{frame}{CAM plants}
    \centering
    \includegraphics[width=0.8\textwidth]{figures/CrassulaPhoto-1.eps}\\
    {\small A species of \species{Crassula}, the genus that originated the name
    crassulacean acid metabolism (CAM).}
\end{frame}

\begin{frame}{CAM ``principle''}
    \begin{itemize}
        \item Initial carboxylation and assimilation are separated
        in time.
        \item At night carboxylation takes place and C from \COtwo is stored as
        malic acid in the vacuole.
        \item[+] The enzyme catalyzing this carboxylation is the
        same as in C4 metabolism: PEP-carboxylase.
        \item During the day decarboxylation takes place, and
        assimilation (including another carboxylation) takes place
        in the Calvin cycle.
    \end{itemize}
\end{frame}

\begin{frame}{CAM mechanism}
    \centering
    \includegraphics[width=\textwidth]{figures/Heldt8.16CAMprinciple.eps}\\
   {\small CAM ``principle'' \autocite[from][]{Heldt1997}.}
\end{frame}

%\section{Energy conversion}
%
%\begin{frame}{Energy conversion}
%    \vspace{-0.07\textheight}
%    \twocolumn[lcolwidth=0.57\linewidth,rcolwidth=0.37\linewidth]{
%    \includegraphics[height=0.75\textheight]{figures/Taiz9.3EnergyConversion.eps}}{%
%    \raggedright
%    \vspace{0.1\textheight}
%    {\small Only about 5\% of the energy is captured in the fixed and reduced carbon
%    \autocite[from][]{TaiZei2006}.}}
%\end{frame}

\section{C3 vs.\ C4 metabolism}

\begin{frame}{History of C4 photosynthesis}
    \begin{itemize}
        \item Nowadays C3 plants account for about 70\% of the world's primary
        productivity.
        \item C4 evolved relatively recently.
        \item 50 million years ago \COtwo concentration in the
        atmosphere ($C_a$) was higher than nowadays.
        \item Much more recently, C4 photosynthesis was more prominent during the glacial periods,
        when $C_a$ was below 200 ppm, than what it is today.
        \item Global increase in \COtwo can be expected to be more
        favourable for C3 than for C4 species.
    \end{itemize}
\end{frame}

\begin{frame}{Use efficiencies}
    \begin{itemize}
        \item C3 photosynthesis is more efficient than C4 photosynthesis in the use of
        light.
        \item C4 photosynthesis is more efficient in the use of
        water and nitrogen.
        \item \COtwo ``pumping'' in C4 requires extra ATP, that
        needs to be produced in the light reactions.
        \item C4 plants do not need to keep stomata so open as C3 plants, and in this way
        they save water.
        \item In C4 plants Rubisco works more efficiently (less
        oxygenation) so less Rubisco is needed to attain the same
        photosynthesis rate. (Less Rubisco means that less nitrogen
        is needed.)
    \end{itemize}
\end{frame}

\begin{frame}{Quantum yield}
    \centering
    \includegraphics[height=0.8\textheight]{figures/Taiz9.9QuantumYieldC3C4.eps}

    {\small Quantum yield in C3 and C4 plants as a function of
    temperature
    \autocite[from][]{TaiZei2006}.}
\end{frame}

\begin{frame}{Conditions}
    \centering
    \includegraphics[height=0.75\textheight]{figures/Taiz9.23FavoredC3C4.eps}

    {\small Conditions favouring C3 and C4 photosynthesis in grasses
    \autocite[from][]{TaiZei2006}.}
\end{frame}

\begin{frame}{Distribution of C3, C4 and CAM}
    \begin{itemize}
        \item Most plants from cool climates have C3 metabolism.
        \item Most shade plants have C3 metabolism.
        \item C4 plants are common in hot and rather dry climates.
        \item In mountains they are more common at the lower
        altitudes where it is warmer.
        \item Many epiphytes are CAM plants.
        \item CAM plants are common in extremely dry and hot
        climates.
        \item Many CAM plants can switch to C3 metabolism when well
        watered.
    \end{itemize}
\end{frame}

\begin{frame}{C4 in mountains}
    \centering
    \includegraphics[width=\textwidth]{figures/Crawley1.26C3C4DistribElevation.eps}\\

    {\small Abundance of C4 grasses along an elevational gradient \autocite[from][]{Crawley1997}.}
\end{frame}

\begin{frame}{Net photosynthesis in C3 vs.\ C4}
    \centering
    \includegraphics[width=0.9\textwidth]{figures/Crawley1.25C3C4Photosyn_gs.eps}\\

    {\small Net photosynthesis and transpiration as a function of leaf conductance
     \autocite[from][]{Crawley1997}.}
\end{frame}

\begin{frame}{Types of metabolism}
    \begin{itemize}
        \item About 70\% of the Earth's primary production is
        from C3 metabolism.
        \item C4 species represent approximately 5\% of all plant
        species, C3 species about 85\%, and CAM species 10\%.
        \item C4 evolved nearly synchronously on different
        continents 5 to 7 million years ago. (At the time \COtwo
        concentration in the atmosphere decreased.)
    \end{itemize}
\end{frame}

\begin{frame}{Types of metabolism}
    \begin{description}
        \item[C3] \species{Triticum aestivum} (wheat), \species{Glycine max} (soybean),
        \species{Pisum sativum} (pea), \species{Betula pendula} (silver birch),
        \species{Pinus sylvestris} (Scots pine).
        \item[C4 NADP-ME] \species{Zea mays} (maize), \species{Sorghum
        sudanense} (Sudan grass), \species{Saccharum officinarum} (sugarcane).
        \item[C4 PCK] \species{Panicum maximum}, \species{Chloris gayana},
        \species{Sporobolus fimbriatus} (grasses).
        \item[C4 NAD-ME] \species{Atriplex spongiosa}, \species{Portulaca
        oleracea} (when well watered), \species{Amaranthus edulis}.
        \item[CAM] \species{Ananas comosus} (pineapple),
        \species{Clusia major} (a tropical tree), \species{Opuntia
        ficus-indica} (nopal, `prickly pear', kaktusviikuna).
        \item[note:] PCK = Phosphoenolpyruvate carboxykinase; ME = Malic enzyme.
        \end{description}
\end{frame}

\begin{frame}{Types of metabolism}
    \begin{description}
        \item[C3/CAM facultative] \species{Mesembryanthemun crystrallinum} (switches to CAM under salinity),
        \species{Clusia minor}.
        \item[C4 NAD-ME/CAM facultative] \species{Portulaca oleracea} (switches to CAM if under drought).
        \item[C3-C4 intermediate] \species{Flaveria floridana}.
        (same genus but C4: \species{F. trinervia}, or C3:
        \species{F. mcdougallii}.)
        \item[CAM in water plants] \species{Isoetes} (fixing of
        \COtwo released by respiration).
    \end{description}
    \small note: CAM and C4 metabolism have appeared several times independently
    during evolution. (CAM also has subtypes: ME-CAM and PCK-CAM.)
\end{frame}

\begin{frame}{Environmental correlations}
    \begin{itemize}
        \item Environmental correlations are just correlations, they
        do not demonstrate on their own cause and effect
        relationships.
        \item In addition to the correlations between high
        temperatures and C4 species, there are correlation patterns for the
        different types of C4 metabolism.
        \item There is negative correlation between higher rainfall and CAM species, for
        tropical epiphytes even at very high rainfall levels.
    \end{itemize}
\end{frame}

\begin{frame}{C4 types and rainfall}
    \centering
    \includegraphics[width=\textwidth]{figures/Lambers43C4grassesAustralia.eps}\\

    {\small Proportion of grasses in Australia with different types
    of photosynthetic metabolism
    \autocite[from][]{LambersEtAl1998}.}
\end{frame}

\begin{frame}{CAM and rainfall}
    \centering
    \includegraphics[height=0.65\textheight]{figures/Lambers50CAMrainfall.eps}\\

    {\small Proportion of epipythic bromeliads that have CAM photosynthetic
    metabolism
    \autocite[from][]{LambersEtAl1998}.}
\end{frame}


%\begin{frame}{CAM mechanism: night}
%%    \vspace{-0.04\textheight}
%    \twocolumn[lcolwidth=0.67\linewidth,rcolwidth=0.27\linewidth]{
%    \includegraphics[height=0.9\textheight]{figures/Heldt8.17CAMnight.eps}}{%
%    \raggedright \vspace{0.8\textheight}
%    {\small CAM mechanism at night \autocite[from][]{Heldt1997}.}}
%\end{frame}
%
%\begin{frame}{CAM mechanism: day}
%%    \vspace{-0.04\textheight}
%    \twocolumn[lcolwidth=0.67\linewidth,rcolwidth=0.27\linewidth]{
%    \includegraphics[height=\textheight]{figures/Heldt8.18CAMday.eps}}{%
%    \raggedright \vspace{0.8\textheight}
%    {\small CAM mechanism during daytime \autocite[from][]{Heldt1997}.}}
%\end{frame}
%
%
%\section{Biochemistry: CAM}
%
%\begin{frame}{CAM plants}
%    \vspace{-0.04\textheight}
%    \twocolumn[lcolwidth=0.57\linewidth,rcolwidth=0.37\linewidth]{
%    \includegraphics[height=0.75\textheight]{figures/CAMphoto-1.eps}}{
%%    \vspace{-0.02\textheight}
%    {\small Some arid habitat plants, typically with CAM.}}
%\end{frame}
%
%\begin{frame}[toc=,bm=]{CAM plants}
%    \vspace{-0.04\textheight}
%    \includegraphics[width=0.8\textwidth]{figures/CrassulaPhoto-1.eps}\\
%%    \vspace{-0.02\textheight}
%    {\small A species of \species{Crassula}, the genus that originated the name
%    crassulacean acid metabolism (CAM).}
%\end{frame}
%
%\begin{frame}{CAM ``principle''}
%    \vspace{-0.04\textheight}
%    \begin{itemize}[type=1]
%        \item<1> Initial carboxylation and assimilation are separated
%        in time.
%        \item<2> At night carboxylation takes place and C from \COtwo is stored as
%        malic acid in the vacuole.
%        \item[+]<2> The enzyme catalyzing this carboxylation is the
%        same as in C4 metabolism: PEP-carboxylase.
%        \item<3> During the day decarboxylation takes place, and
%        assimilation (including another carboxylation) takes place
%        in the Calvin cycle.
%    \end{itemize}
%\end{frame}
%
%\begin{frame}[toc=,bm=]{CAM mechanism}
%    \vspace{-0.06\textheight}
%    \includegraphics[width=\textwidth]{figures/Heldt8.16CAMprinciple.eps}\\
%    \vspace{-0.02\textheight}
%    {\small CAM ``principle'' \autocite[from][]{Heldt1997}.}
%\end{frame}
%
%\begin{frame}{CAM mechanism: night}
%%    \vspace{-0.04\textheight}
%    \twocolumn[lcolwidth=0.67\linewidth,rcolwidth=0.27\linewidth]{
%    \includegraphics[height=0.9\textheight]{figures/Heldt8.17CAMnight.eps}}{%
%    \raggedright \vspace{0.8\textheight}
%    {\small CAM mechanism at night \autocite[from][]{Heldt1997}.}}
%\end{frame}
%
%\begin{frame}{CAM mechanism: day}
%%    \vspace{-0.04\textheight}
%    \twocolumn[lcolwidth=0.67\linewidth,rcolwidth=0.27\linewidth]{
%    \includegraphics[height=\textheight]{figures/Heldt8.18CAMday.eps}}{%
%    \raggedright \vspace{0.8\textheight}
%    {\small CAM mechanism during daytime \autocite[from][]{Heldt1997}.}}
%\end{frame
%}
%
%

  \section*{References}
  \begin{frame}[t,allowframebreaks]
    \frametitle{References}
    \nocite{Heldt2004}
    \nocite{Lambers2008}
    \printbibliography
  \end{frame}

\end{document}
